\chapter{Diagrama de flujo}
  
\begin{center}
  \begin{tikzpicture}[node distance=.8cm]  
    \node (start) [startstop] {Inicio};
    \node (process1) [process, below = of start] {Revisión bibliográfica de los modelos de SLM};
    \node (process2) [process, below =of process1] {Análisis de las ecuaciones adaptándolas a la configuración del dispositivo};
    \node (process3) [process, below=of process2] {Simulación del sistema no lineal en los puntos de equilibrio};
    \node (process4) [process, below=of process3] {Linealización alrededor de los puntos de equilibrio.};
    \node (process5) [process, below=of process4] {Diseño y simulación del controlador lineal que estabilice, sin efecto atascamiento-deslizamiento,  utilizando un controlador lineal con doble efecto integral.};
    \node (process6) [process, below=of process5] {Pruebas al controlador incluyendo el atascamiento-deslizamiento en el modelo lineal.};
    \node (process7) [process, below=of process6] {Pruebas del controlador en el modelo no lineal sin efecto atascamiento-deslizamiento};
    \node (process8) [process, below=of process7] {Pruebas del controlador en el modelo no lineal con efecto atascamiento-deslizamiento.};
    \node (end) [startstop, below = of process8] {Final};
    
    
    \draw[->,thick,>=stealth] (start) --  (process1); 
    \draw[->,thick,>=stealth] (process1) -- (process2);
    \draw[->, thick,>=stealth] (process2)--(process3);
    \draw[->, thick,>=stealth] (process3)--(process4);
    \draw[->, thick,>=stealth] (process4)--(process5);
    \draw[->, thick,>=stealth] (process5)--(process6);
    \draw[->, thick,>=stealth] (process6)--(process7);
    \draw[->, thick,>=stealth] (process7)--(process8);
    \draw[->,thick,>=stealth] (process8) -- (end);

  \end{tikzpicture}
\end{center}